\begin{statementbox}

	\textbf{\textcolor{CStatement}{对象与约定}}

	给定 $n$ 个实数 $x_1,x_2,\dots,x_n$，其中 $n\in\mathbb N^*$。将它们从小到大排列为
	$x_{(1)}\le x_{(2)}\le\cdots\le x_{(n)}$。本结论中的“方差”采用高中描述统计口径，分母为 $n$；推断统计中分母为 $n-1$ 的无偏样本方差不在此处讨论。

	\medskip
	\textbf{\textcolor{CStatement}{集中趋势}}
	\begin{description}[style=nextline,leftmargin=2.8em,labelsep=0.5em,itemsep=0.45em]
		\item[\textbf{平均数}]
			\[
				\bar x=\frac1n\sum_{i=1}^{n}x_i.
			\]
			它使用每一个数据值，表示这组数据的平衡水平。

		\item[\textbf{中位数}]
			\[
				\operatorname{Med}=\begin{cases}
					x_{((n+1)/2)},&n\text{ 为奇数},\\[3pt]
					\dfrac{x_{(n/2)}+x_{(n/2+1)}}2,&n\text{ 为偶数}.
				\end{cases}
			\]
			它只由排序后的中间位置决定。

		\item[\textbf{众数}]
			出现次数达到最大值的数据。众数可能不唯一；若所有数据互不相同，高中统计中通常称“没有众数”。
	\end{description}

	\medskip
	\textbf{\textcolor{CStatement}{离散程度}}
	\[
		s^2=\frac1n\sum_{i=1}^{n}(x_i-\bar x)^2
		=\frac1n\sum_{i=1}^{n}x_i^2-\bar x^2,
		\qquad s=\sqrt{s^2}.
	\]
	恒有 $s^2\ge0$；并且
	\[
		s^2=0\iff x_1=x_2=\cdots=x_n.
	\]
	方差的量纲是原数据量纲的平方，标准差与原数据同量纲。用方差或标准差比较波动时，两组数据应描述同一变量并采用相同单位。

	\medskip
	\textbf{\textcolor{CStatement}{频数表形式}}

	若不同取值 $a_1,\dots,a_k$ 的频数分别为非负整数 $f_1,\dots,f_k$，总频数 $F=\sum_{j=1}^{k}f_j>0$，则
	\[
		\bar x=\frac{\sum_{j=1}^{k}f_ja_j}{F},
		\qquad
		s^2=\frac{\sum_{j=1}^{k}f_j(a_j-\bar x)^2}{F}.
	\]
	这不是另一套定义：它只是把每个 $a_j$ 按频数 $f_j$ 重复写出的简记。

\end{statementbox}
